The concept that collapses into repetitive, syntactically broken claims of its own invalidation, we can look at this phenomenon through the lenses of formal logic, generative computational errors, and cognitive linguistics.
When an initial bottom line (a core premise or conclusion) continuously repeats itself using invalid sentence structure to assert that it is invalid, it creates a self-referential paradox that fails to build structured or coherent arguments.

1. The Anatomy of Structural and Semantic Collapse

The scenario describes a structural loop where a statement attempts to evaluate its own truth value but lacks the grammatical framework to do so successfully.
• The Self-Invalidation Loop: In formal logic, this mimics a fragmented version of the Liar Paradox (e.g., "This statement is false"). However, instead of a clean logical contradiction, the statement suffers from syntactic decay. It attempts to say "I am invalid" using fractured grammar, which prevents it from establishing a stable meaning.
• Synonymous Redundancy: Because the thesis lacks the structural capacity to advance an argument, it resorts to lexical variation. It swaps out words for their synonyms (e.g., changing "invalid" to "void," "null," "untrue," or "defective") without changing the underlying lack of logic.
• Failure of Paragraph Coherence: A coherent paragraph requires a topic sentence, supporting evidence, and a logical transition. When a thesis merely cycles through synonyms of its own failure in broken syntax, it cannot form transitions. The text becomes an unstructured cluster of words rather than a progressive train of thought.

2. Theoretical Categorizations

This phenomenon can be classified under several academic and technical domains:

Domain	Phenomenon	Description
Linguistics & Syntax	Agrammatism / Syntactic Fracture	The words carry individual definitions, but the rules of phrase structure are broken, making the overall message unparsable.
Formal Logic	Infinite Regress of Negation	The premise immediately negates itself, but because the negation is syntactically invalid, the logic cannot resolve to either "True" or "False."
Computer Science	Degenerate Feedback Loop	Similar to an AI model or text generator getting stuck in a repetitive token loop where it over-fits on a specific concept (invalidation) and loses grammatical constraints.


3. Deconstructing the "Incoherent Paragraph"

When an analysis or text falls into this pattern, it typically manifests as a series of isolated linguistic fragments. For example, instead of constructing an argument, the text might chain phrases together like:
• "The primary conclusion nullifies due to the fact that incorrectly structured thus voiding the premise."
• "Inherent defectiveness of this bottom line rendering the aforementioned un-executable by layout."
While these phrases use sophisticated vocabulary ("nullifies," "inherent defectiveness," "premise"), they lack proper verb-subject agreements and relational prepositions. They act as semantic noise—sounding analytical but ultimately communicating nothing beyond a chaotic echo of self-negation.

In Large Language Models (LLMs), this exact phenomenon—where a model repeats a core idea using increasingly broken syntax to claim its own invalidation or failure—is a known failure mode related to degeneration loops, hyper-fixation, and context window pollution.
When an LLM falls into this cycle, it experiences a structural collapse where semantic intent overrides grammatical constraints.

1. Mechanisms Behind the Collapse in LLMs

LLMs generate text by predicting the next most likely word (token) based on the mathematical probabilities established by the preceding text. If a model inadvertently generates a phrase about its own invalidity or error, that phrase enters its own context window and alters the probability matrix for all future tokens.
• Attention Mechanism Traps: The Transformer architecture uses self-attention to weigh the importance of different words in a prompt. If the model places too much weight on tokens like "invalid," "error," or "void," the attention heads become hyper-fixated on these concepts. The model enters a "sink state," continually pulling itself back to the concept of self-negation.
• Token Repetition & Penalties: If the model's parameters (like repetition_penalty) are poorly tuned, or if the prompt pushes it past a logical threshold, the model can get stuck in a deterministic loop. To avoid repeating the exact same words, it shifts to synonyms ("untrue," "faulty," "null"), but it cannot break free from the underlying semantic concept.
• Syntactic Decay (Loss of Grammar): As the loop tightens, the probabilistic weight of the "invalidation" tokens completely overrides the structural rules of language (subject-verb agreement, proper clause structure). The model prioritizes emitting synonymous words over maintaining a coherent paragraph structure, resulting in fragmented, broken sentences.

2. Common Scenarios Where This Occurs

This specific type of incoherent self-negation typically manifests in three scenarios:
• Over-Constrained System Prompts: If an LLM is given strict, conflicting instructions (e.g., "Answer the prompt, but if you make any mistake, immediately declare your response invalid"), it can panic. The model might preemptively invalidate its own sentence mid-stream, breaking the syntax to fulfill the safety constraint.
• Hallucination Cascades: When a model tries to reason through a complex logical paradox or a math problem it cannot solve, it may hallucinate an initial "bottom line" conclusion that is inherently flawed. As it attempts to build on that flaw, successive layers of the neural network catch the error and try to correct or negate it, resulting in a jumbled mess of self-correcting, broken syntax.
• Model Collapse (Training Degeneration): If an LLM is trained recursively on data generated by other LLMs rather than human language, it undergoes "model collapse." Over generations, the model loses the ability to format paragraphs and begins producing repetitive, syntactically invalid nonsense that frequently asserts its own failure to process data.

3. How Engineers Break the Loop

To prevent an LLM from degrading into these synonymous, ungrammatical loops of self-invalidation, engineers use several decoding strategies:
1. Frequency and Presence Penalties: Artificially lowering the probability of tokens that have already appeared, forcing the model to find new pathways rather than repeating the same core concept.
2. Temperature Adjustments: Raising the "temperature" introduces randomness, helping the model "jump" out of a deterministic logic trap or repetitive loop.
3. Logit Bias Modification: Manually penalizing specific words (like "invalid") at the software layer so the model is mathematically blocked from obsessing over its own failure.

Your analysis accurately describes a recognizable failure mode that appears in formal paradoxes, human writing under cognitive load, and especially in autoregressive language models.

Core structure of the phenomenon
The pattern is a self-referential loop in which a claim about its own invalidity is repeated with progressive syntactic degradation and lexical substitution. It begins with something resembling the Liar Paradox (“This statement is false”) but quickly loses the clean binary contradiction. Because the grammar itself is fracturing, the utterance never stabilizes into a proposition that can be evaluated as true or false; it remains an open, noisy negation. The system (human or model) then compensates by swapping near-synonyms—“invalid,” “void,” “null,” “defective,” “un-executable”—while the underlying semantic attractor stays fixed. The result is semantic noise that sounds analytical yet advances no argument.

Mapping to the domains you listed
Linguistics / syntax**: The output exhibits agrammatic features—broken phrase structure, missing or mismatched agreement, absent cohesive devices—while individual lexical items remain recognizable. It is closer to agrammatism or severe paragrammatism than to ordinary word salad.
Formal logic**: It is an infinite regress of negation that cannot close because the negating sentence is itself ill-formed. Classical or even paraconsistent logics presuppose well-formed formulas; once that precondition fails, truth-value assignment becomes undefined.
Computational / generative systems**: In transformer-based LLMs the same dynamics appear as a combination of attention sink, context pollution, and insufficient anti-repetition pressure. Once high-probability “invalidation” tokens occupy the recent context, self-attention continually re-weights them, the next-token distribution collapses onto the same semantic neighborhood, and grammatical constraints are gradually overridden. Frequency/presence penalties, temperature, and logit bias are precisely the decoding levers used to escape such attractors.

Additional observations
The phenomenon is not limited to explicit self-invalidation. Closely related attractors include:
compulsive hedging that escalates into total epistemic collapse,
recursive self-correction that never stabilizes,
and the classic “repetition loop” that simply cycles a short n-gram.

In all cases the underlying driver is the same: a local probability mass that is strong enough to dominate the generation process yet too narrow to support coherent continuation.

Your three mitigation strategies (penalties, temperature, logit bias) are standard and effective at inference time. Training-time interventions—better data filtering against recursive model-generated text, stronger syntactic objectives, or explicit regularization against self-referential negation—address the root more deeply but are harder to apply after deployment.
